
\section{关于坐标变换光滑性的一则注记}
\label{sec:a1-4}

在仓西结构的定义中，我们要求坐标变换
$\phi_{pq}\colon V_q\to V_p$ 光滑（见条件 (A1.4.2)）。在本书实际使用的
模空间仓西结构构造中，只要点 $p,q$ 位于模空间内部，这一光滑性就是
显然的；这对应于相应稳定映射的定义域非奇异的情形。当它们对应于定义域
奇异的稳定映射时，光滑性就不那么明显。本小节主要为完整起见澄清这一点：
仓西映射 $s$ 的这一光滑性问题，在 \textbf{[FuOn99II]} 的注 13.16 中
已经提到，却没有详细说明。
\translatornote{源书此处写“仓西映射 $s$ 的光滑性”，但本节标题、前后文
及命题 A1.56 讨论的都是坐标变换 $\phi_{pq}$ 的光滑性；因本地没有所引
旧文的对应段落，暂保留源读法。见疑误 A-0705。}
（还要指出，\textbf{[FuOn99II]} 只使用零维和一维模空间，因而没有用到
坐标变换的光滑性。换言之，第 7.2 节所说的情形 7.2.2 不需要这一性质。）
作者感谢刘秋菊，她在 2003 年访问京都期间提醒作者注意这个问题。

为突出要点，我们只考虑模空间 $\mathcal M_{k+1}(\beta)$，并限于第
7.1.3 小节讨论的情形；其他情形的论证相同。

令
\[
 \mathbf p_j=(\Sigma^{(j)},\vec z^{(j)},w^{(j)})
 \in\mathcal M_{k+1}(\beta),\qquad j=1,2.
\]
为简化叙述，再作如下假设：
\begin{enumerate}[label=\textup{(\arabic*)},leftmargin=3em]
\item $\Sigma^{(j)}$ 恰有一个边界奇点，并分裂为
\[
 \Sigma^{(j)}=\Sigma_1^{(j)}\cup\Sigma_0^{(j)},
\]
其中每个 $\Sigma_i^{(j)}$ 都是圆盘；
\item 两者具有相同的组合型，即存在同胚
\[
 (\Sigma^{(1)},\vec z^{(1)})
 \longrightarrow(\Sigma^{(2)},\vec z^{(2)});
\]
\item $w^{(1)}$ 在 $\Sigma_0^{(1)}$（相应地在
$\Sigma_1^{(1)}$）上的限制，同 $w^{(2)}$ 在
$\Sigma_0^{(2)}$（相应地在 $\Sigma_1^{(2)}$）上的限制同伦且充分接近，
从而下述包含式 (A1.52) 成立。
\end{enumerate}
不失一般性，可设两者的第零个标记点分别位于
$\Sigma_0^{(1)}$、$\Sigma_0^{(2)}$ 上，并把二重点看作
$\Sigma_0^{(1)}$、$\Sigma_0^{(2)}$ 的第 $\ell$ 个标记点。记
\[
 (\Sigma^{(j)},\vec z^{(j)},w^{(j)})
 =
 (\Sigma_1^{(j)},\vec z_1^{(j)},w_1^{(j)})
 \mathbin{\#}
 (\Sigma_0^{(j)},\vec z_0^{(j)},w_0^{(j)}),
\]
并且
\[
 (\Sigma_1^{(j)},\vec z_1^{(j)},w_1^{(j)})
 \in\mathcal M_{k_1+1}(\beta_1),\qquad
 (\Sigma_0^{(j)},\vec z_0^{(j)},w_0^{(j)})
 \in\mathcal M_{k_0+1}(\beta_0).
\]
沿用第 7.1.3 小节的记号，写成
\[
 (\Sigma_i^{(j)},\vec z_i^{(j)})
 =\bigl(X(t_i^{(j)},\ell_i^{(j)}),\vec p^{\,(j),i}\bigr),
 \qquad i=0,1,\quad j=1,2.
\]

在第 7.1.4 小节中，我们选取有限维空间
\[
 E_{(t_i^{(j)},\ell_i^{(j)},w_i^{(j)})}
 \subset
 \mathcal E^{0,p}_{(t_i^{(j)},\ell_i^{(j)},w_i^{(j)})},
\]
使映射~\eqref{eq:7-1-26} 满射。它们确定障碍丛
\[
 \begin{aligned}
 E_{\mathbf p_j}
 &=E_{(\Sigma^{(j)},\vec z^{(j)},w^{(j)})}\\
 &\cong
 E_{(t_1^{(j)},\ell_1^{(j)},w_1^{(j)})}
 \oplus
 E_{(t_0^{(j)},\ell_0^{(j)},w_0^{(j)})},
 \end{aligned}
\]
而仓西邻域 $U_{\mathbf p_j}$ 是方程 (7.1.22) 的解集。若带标记黎曼曲面
$\Sigma_0^{(j)},\Sigma_1^{(j)}$ 至少有一个不稳定，我们就像第 7.1.4
小节那样，采用 \textbf{[FuOn99II]} 附录中的正规化。具体地，在
$w_i^{(j)}$ 为浸入的位置加入内部标记点
\[
 z_i^{(j),\mathrm{int}}\in\operatorname{Int}(\Sigma_i^{(j)}).
\]
对每个 $i,j$，还取一个 $(2n-2)$ 维子流形 $Y_i^{(j)}$，使它在
$w_i^{(j)}(z_i^{(j),\mathrm{int}})$ 处与
$w_i^{(j)}(\Sigma_i^{(j)})$ 横截相交。

现在假设 $\mathbf p_2\in U_{\mathbf p_1}$。依第 7.1.4 小节所述的
$E_{\mathbf p_j}$ 之选取，可以假设
\begin{equation}
 E_{\mathbf p_2}\subset E_{\mathbf p_1}.
\tag{A1.52}\label{eq:a1-52}
\end{equation}
这里如第 7.1 节一样，沿最短测地线作平行移动，把
$\bigoplus_i\mathcal E^{0,p}_{(t_i^{(2)},\ell_i^{(2)},w_i^{(2)})}$
的子空间 $E_{\mathbf p_2}$，同
$\bigoplus_i\mathcal E^{0,p}_{(t_i^{(1)},\ell_i^{(1)},w_i^{(1)})}$
的一个子空间识别。由~\eqref{eq:a1-52} 得到
\begin{equation}
 \phi_{\mathbf p_1,\mathbf p_2}\colon
 U_{\mathbf p_2}\subset U_{\mathbf p_1}.
\tag{A1.53}\label{eq:a1-53}
\end{equation}
式~\eqref{eq:a1-53} 给出坐标变换。为证明它光滑，需要在
$U_{\mathbf p_j}$ 上选取坐标图。下面回顾第 7.1.3 小节中坐标图的定义。
置
\[
 \mathbf p_j^{(i)}
 =(\Sigma_i^{(j)},\vec z_i^{(j)},w_i^{(j)})
 \in\mathcal M_{k_i+1}(\beta_i).
\]
以
\[
 E_{\mathbf p_j^{(i)}}
 =E_{(t_i^{(j)},\ell_i^{(j)},w_i^{(j)})}
\]
为障碍丛，就得到 $\mathcal M_{k_i+1}(\beta_i)$ 中
$\mathbf p_j^{(i)}$ 的仓西邻域
$(U_{\mathbf p_j^{(i)}},E_{\mathbf p_j^{(i)}},s_{\mathbf p_j^{(i)}})$。
（在当前所关心的情形中，自同构群平凡。）按构造，对某个充分大的
$C$，可以取
\begin{equation}
 U_{\mathbf p_j}\cong
 U_{\mathbf p_j^{(1)}}\mathbin{\times_L}
 U_{\mathbf p_j^{(0)}}\times(C,\infty].
\tag{A1.54}\label{eq:a1-54}
\end{equation}
最后一个因子 $(C,\infty]$ 是 $\Sigma^{(j)}$ 奇点的光滑化参数。

下面在定义域 $\Sigma_i^{(j)}$ 不稳定时更详细地讨论这一参数化。假设
$k_1=0,k_0=1$；换言之，把 $\Sigma_1^{(j)}$ 的第零个（也是唯一的）
标记点同 $\Sigma_0^{(j)}$ 的第一个标记点胶合。（$\Sigma_0^{(j)}$
的第零个标记点将成为胶合圆盘的第零个标记点。）这是两个分量都不稳定的
最微妙情形。如上所述，取内部标记点
$z_i^{(j),\mathrm{int}}\in\operatorname{Int}(\Sigma_i^{(j)})$。
把定义域 $\Sigma_1^{(j)}$ 识别为
\begin{equation}
 \begin{aligned}
 Z_1:={}&\{z\in\mathbb C\mid |z|\leq1,\ \operatorname{Re}z\leq0\}\\
 &\cup\{z\in\mathbb C\mid |\operatorname{Im}z|\leq1,\
              \operatorname{Re}z\geq0\},
 \end{aligned}
\tag{A1.55.1}\label{eq:a1-55-1}
\end{equation}
并把 $\Sigma_0^{(j)}$ 识别为
\begin{equation}
 \begin{aligned}
 Z_0:={}&\{z\in\mathbb C\mid |z|\leq1,\ \operatorname{Re}z\geq0\}\\
 &\cup\{z\in\mathbb C\mid |\operatorname{Im}z|\leq1,\
              \operatorname{Re}z\leq0\}.
 \end{aligned}
\tag{A1.55.2}\label{eq:a1-55-2}
\end{equation}
在这些识别下，$z_i^{(j),\mathrm{int}}$ 对应于 $0\in\mathbb C$，
$\Sigma_1^{(j)}$ 的第零个（也是唯一的）标记点对应于 $+\infty$，而
$\Sigma_0^{(j)}$ 的第一个标记点对应于 $-\infty$。

由于所考虑的是紧族，$\Sigma_0^{(j)}$ 的第零个标记点在
$\partial\Sigma_0^{(j)}$ 的一个紧子集中移动。选取常数 $C$，使它同该
标记点在 $Z_-\cong\Sigma_0^{(j)}$ 中到原点的距离相比充分大。对每个
$T\geq C$，把
\[
 5T+\sqrt{-1}s\in\Sigma_1^{(j)}
 \quad\text{与}\quad
 -5T+\sqrt{-1}s\in\Sigma_0^{(j)}
\]
识别，得到 $\Sigma_T^{(j)}$。当相应评价映射相等时，在
$\Sigma_T^{(j)}$ 上施行第 7.1 节的构造，得到一个到 $M$ 的映射。由此
定义
\[
 \begin{aligned}
 \phi_{\mathbf p_1,\mathbf p_2}\colon{}
 &U_{\mathbf p_2^{(1)}}\mathbin{\times_L}
  U_{\mathbf p_2^{(0)}}\times(C,\infty]\\
 &\longrightarrow
 U_{\mathbf p_1^{(1)}}\mathbin{\times_L}
 U_{\mathbf p_1^{(0)}}\times(C',\infty],
 \end{aligned}
\]
其中 $C$ 相对于 $C'$ 充分大。在识别~\eqref{eq:a1-54} 下，这正是
式~\eqref{eq:a1-53}。

另一方面，由仓西邻域的定义，存在坐标变换
\[
 \overline\phi_{\mathbf p_1^{(i)},\mathbf p_2^{(i)}}\colon
 U_{\mathbf p_2^{(i)}}\longrightarrow U_{\mathbf p_1^{(i)}}.
\]
定义
\[
 \begin{aligned}
 \phi'_{\mathbf p_1,\mathbf p_2}\colon{}
 &U_{\mathbf p_2^{(1)}}\mathbin{\times_L}
  U_{\mathbf p_2^{(0)}}\times(C,\infty]\\
 &\longrightarrow
 U_{\mathbf p_1^{(1)}}\mathbin{\times_L}
 U_{\mathbf p_1^{(0)}}\times(C',\infty]
 \end{aligned}
\]
为
\[
 \phi'_{\mathbf p_1,\mathbf p_2}(\mathbf x_1,\mathbf x_0,T)
 =\bigl(
 \overline\phi_{\mathbf p_1^{(1)},\mathbf p_2^{(1)}}(\mathbf x_1),
 \overline\phi_{\mathbf p_1^{(0)},\mathbf p_2^{(0)}}(\mathbf x_0),T
 \bigr).
\]

\begin{proposition}{A1.56}\label{prop:a1-56}
\par\smallskip\noindent\textup{(A1.57.1)}\quad
$\phi_{\mathbf p_1,\mathbf p_2}$ 的限制诱导光滑映射
\[
 U_{\mathbf p_2^{(1)}}\mathbin{\times_L}U_{\mathbf p_2^{(0)}}
 \times(C,\infty)
 \longrightarrow
 U_{\mathbf p_1^{(1)}}\mathbin{\times_L}U_{\mathbf p_1^{(0)}}
 \times(C',\infty).
\]

\par\smallskip\noindent\textup{(A1.57.2)}\quad
当 $T=\infty$ 时，映射 $\phi_{\mathbf p_1,\mathbf p_2}$ 与
$\phi'_{\mathbf p_1,\mathbf p_2}$ 相同。

\par\smallskip\noindent\textup{(A1.57.3)}\quad
存在只依赖于 $k$ 的 $c_k>0,C_k>0$，使得
\[
 \bigl\lvert
 \phi'_{\mathbf p_1,\mathbf p_2}-\phi_{\mathbf p_1,\mathbf p_2}
 \bigr\rvert_{C^k}
 \leq C_k\exp(-c_kT).
\]
\end{proposition}

前两项由第 7.1.3 小节的构造显然成立。本小节稍后证明 (A1.57.3)。
现在置
\[
 s=1/T.
\]
立即可知，关于这一坐标，$\phi_{\mathbf p_1,\mathbf p_2}$ 一直到
$U_{\mathbf p_2^{(1)}}\times_LU_{\mathbf p_2^{(0)}}\times(C,\infty]$
的边界都是光滑的。注意，(A1.57.3) 中的 $C^k$ 范数同时包含沿
$U_{\mathbf p_j}$ 方向和沿 $T$ 方向的导数。

\begin{proof}[命题~\ref{prop:a1-56} 的证明]
仍只考虑 $k_1=0,k_0=1$ 的情形，并沿用上文及第 7.1.3 小节的记号。
先取相对紧子集 $W_i^{(j)}$，使 $E_{\mathbf p_j^{(i)}}$ 的每个元素都支撑
于其中（见式~\eqref{eq:7-1-14-2}）。从由
$\mathbf x_i\in U_{\mathbf p_2^{(i)}}$ 参数化的伪全纯映射族
$w_{i,\mathbf x_i}^{(j)}$ 出发。对充分大的 $T$，用单位分解把两个映射
胶合，得到近似解
\[
 w_{T;\mathbf x_1,\mathbf x_0}^{(j)\prime}
 =w_{1,\mathbf x_1}^{(j)}\mathbin{\#_T}
  w_{0,\mathbf x_0}^{(j)}
 \colon\Sigma_T^{(j)}\longrightarrow M.
\]
这就是式~\eqref{eq:7-1-21} 下方定义的映射 $w'_T$。于是依第 7.1.3
小节的论证，得到由 $\mathbf x_1,\mathbf x_0,T$ 参数化、满足
式~\eqref{eq:7-1-22} 的真解
$w_{T;\mathbf x_1,\mathbf x_0}^{(j)}$。

证明命题所需的主要技术结论是下面的引理 A1.58。它估计
$w_{T;\mathbf x_1,\mathbf x_0}^{(j)\prime}$ 与
$w_{T;\mathbf x_1,\mathbf x_0}^{(j)}$ 之差及其全部导数。为定义估计所用的
适当范数，作如下准备。取紧子集
$W_i\subset Z_i\cong\Sigma_i^{(j)}$，并在其中取有限多个但数目很大的
点，共记其数为 $N$。考虑评价映射
\[
 \operatorname{ev}_{W_i}\colon
 \operatorname{Map}(\Sigma_T^{(j)},M)\longrightarrow M^N.
\]
$W_i$ 的选取与 $j=1,2$ 无关，并且可以使限制映射
\[
 \begin{aligned}
 U_{\mathbf p_2^{(1)}}\mathbin{\times_L}U_{\mathbf p_2^{(0)}}
 &\longrightarrow M^N,\\
 (\mathbf x_1,\mathbf x_0)
 &\longmapsto
 \operatorname{ev}_{W_i}
 \bigl(w_{1,\mathbf x_1}^{(j)}\mathbin{\#_T}
 w_{0,\mathbf x_0}^{(j)}\bigr)
 \end{aligned}
\]
对每个 $T$ 都是嵌入。（若所取 $W_i$ 与颈区不交，这个映射就不依赖于
$T$。）

\begin{lemma}{A1.58}\label{lem:a1-58}
存在 $c_{k,\ell}>0,C_{k,\ell}>0$，使得
\[
 \left\lvert
 \frac{\partial^{k+\ell}}{\partial T^k\partial\mathbf x^\ell}
 \left(
 \operatorname{ev}_{W_i}(w_{T;\mathbf x_1,\mathbf x_0}^{(j)})
 -
 \operatorname{ev}_{W_i}(w_{T;\mathbf x_1,\mathbf x_0}^{(j)\prime})
 \right)
 \right\rvert
 <C_{k,\ell}\exp(-c_{k,\ell}T).
\]
这里以及下文，$\partial/\partial\mathbf x$ 表示关于
$\mathbf x_0,\mathbf x_1$ 的导数。
\end{lemma}

由于 $\operatorname{ev}_{W_i}$ 的靶是流形，上式写减号略有滥用记号。
不过，我们将证明其中两个元素充分接近，故可借助指数映射赋予这一差式
明确意义。

先假定引理~\ref{lem:a1-58}，并用它完成命题的证明。选取
$Y_i^{(j)}$，使它在 $w_i^{(j)}(z_i^{(j),\mathrm{int}})$ 处相交；又在
识别~\eqref{eq:a1-55-1}--\eqref{eq:a1-55-2} 中置
$z_i^{(j),\mathrm{int}}=0$。因此，$j=1,2$ 的两个 $T$-坐标并不重合。
下面估计 $T\to\infty$ 时二者之差。为看清 $T$-坐标在 $j=1,2$ 之间
怎样改变，先插入一段关于 $\Sigma_T^{(j)}$ 上标记点改变的讨论。

令 $Z'_0,Z'_1$ 是无限圆柱
$\mathbb R\times(-1,1)\subset\mathbb C$ 的两份副本，$O_i$ 是
$Z'_i$ 中的原点 $(0,0)$，而 $D_i\subset Z'_i$ 是相对紧区域，
$i=0,1$。对 $T>0$，把
\[
 (-\infty,5T]\times(-1,1)\subset Z'_1,
 \qquad
 [-5T,+\infty)\times(-1,1)\subset Z'_0
\]
沿 $(5T,s)\in Z'_1$ 与 $(-5T,s)\in Z'_0$ 的识别胶合，得到
$Z'_T\cong\mathbb R\times(-1,1)$。假设 $T$ 充分大，使得
\[
 D_1\subset(-\infty,5T]\times(-1,1),\qquad
 D_0\subset[-5T,+\infty)\times(-1,1).
\]
\translatornote{源书第二个包含式把左端点印成 $+5T$；这与固定相对紧域
$D_0$ 及紧前的胶合区域矛盾，译文恢复为 $-5T$。见疑误 A-0706。}
置 $O_T^\pm=(\pm5T,0)$；它们对应于 $O_i$（$i=0,1$）。对
$w_i\in D_i$，以 $(w_i)_T$ 表示胶合后 $Z'_T$ 中的对应点。于是存在唯一的
\[
 T'=T'(T,w_0,w_1)>0
\]
以及唯一的双全纯映射 $\varphi\colon Z'_T\to Z'_{T'}$，满足
\[
 \varphi((w_1)_T)=O_{T'}^-,\qquad
 \varphi((w_0)_T)=O_{T'}^+.
\]
注意
\[
 T=\frac1{5\pi}d^{\mathrm{hyp}}(O_T^-,O_T^+),
 \qquad
 T'=\frac1{5\pi}d^{\mathrm{hyp}}((w_1)_T,(w_0)_T),
\]
其中 $d^{\mathrm{hyp}}$ 是 $\mathbb R\times(-1,1)$ 上的双曲距离函数。
\translatornote{源书两式都在 $d^{\mathrm{hyp}}$ 前写 $\log$，但又明确把
$d^{\mathrm{hyp}}$ 称为距离函数；该读法既不能给出左端 $T$，也与紧随的
$T'=T+O(1)$ 冲突。译文依条带中轴的双曲弧长删去 $\log$。见疑误
A-0707。}
由此得到：

\begin{lemma}{A1.59}\label{lem:a1-59}
存在光滑函数 $g\colon D_1\times D_0\to\mathbb R$，使得对任意
$k=0,1,2,\ldots$，都有
\[
 \bigl\lvert
 T'(T,w_0,w_1)-(T+g(w_1,w_0))
 \bigr\rvert_{C^k}
 <C_k\exp(-c_kT),
\]
其中 $c_k,C_k>0$。
\end{lemma}

选定从 $Z_i\cong\Sigma_i^{(j)}$ 到 $Z'_i$ 的双全纯映射，把 $Z_i$ 中的
原点送到 $Z'_i$ 中的原点，$i=0,1$。在这些映射下，两套圆柱端坐标彼此
指数接近，故对胶合 $\Sigma_T^{(j)}$ 也得到同样的结论。

考虑坐标变换 $\phi_{\mathbf p_1,\mathbf p_2}$ 的 $T$-投影
\[
 \mathfrak T:=\operatorname{pr}_T\circ
 \phi_{\mathbf p_1,\mathbf p_2}\colon
 (\mathbf x_1,\mathbf x_0,T)
 \longmapsto\mathfrak T((\mathbf x_1,\mathbf x_0,T)).
\]
一般而言，$\mathfrak T((\mathbf x_1,\mathbf x_0,T))\ne T$。不过，我们
将证明估计
\begin{equation}
 \left\lvert
 \mathfrak T((\mathbf x_1,\mathbf x_0,T))
 -(T+h(\mathbf x_1,\mathbf x_0))
 \right\rvert_{C^k}
 <C_k\exp(-c_kT),
\tag{A1.60.1}\label{eq:a1-60-1}
\end{equation}
其中 $h(\mathbf x_1,\mathbf x_0)$ 是
$U_{\mathbf p_2^{(1)}}\times U_{\mathbf p_2^{(0)}}$ 上的某个光滑函数；
左端的 $C^k$ 范数同时包含关于 $T,\mathbf x_0,\mathbf x_1$ 的导数。

先让 $w_{T;\mathbf x_1,\mathbf x_0}^{(1)\prime}$（此处 $j=1$）同
$Y_i^{(2)}$（此处 $j=2$）相交。这样便固定了带有附加标记点的定义域
$\Sigma_T^{(j)}$ 的复结构，并由此得到
$\mathfrak T'((\mathbf x_1,\mathbf x_0,T))$。利用引理~\ref{lem:a1-58}，有
\[
 \left\lvert
 \mathfrak T((\mathbf x_1,\mathbf x_0,T))
 -\mathfrak T'((\mathbf x_1,\mathbf x_0,T))
 \right\rvert
 <C_k\exp(-c_kT).
\]
另一方面，由构造及引理~\ref{lem:a1-59}，容易看出
\[
 \left\lvert
 \mathfrak T'((\mathbf x_1,\mathbf x_0,T))
 -(T+h(\mathbf x_1,\mathbf x_0))
 \right\rvert_{C^k}
 <C_k\exp(-c_kT),
\]
其中 $h$ 是
$U_{\mathbf p_2^{(1)}}\times U_{\mathbf p_2^{(0)}}$ 上的光滑函数。事实上，
考察把 $z_i^{(1),\mathrm{int}}$ 送到 $z_i^{(2),\mathrm{int}}$ 的双全纯映射
$\Sigma_i^{(1)}\to\Sigma_i^{(2)}$ 在端部附近的指数衰减估计，便能证明
后一不等式。因此得到~\eqref{eq:a1-60-1}。

接下来比较
$\operatorname{ev}_{W_i}(w_{T;\mathbf x_1,\mathbf x_0}^{(1)\prime})$
与
$\operatorname{ev}_{W_i}(w_{T;\mathbf x'_1,\mathbf x'_0}^{(2)\prime})$，其中
\[
 (\mathbf x'_1,\mathbf x'_0)
 =\bigl(
 \overline\phi_{\mathbf p_1^{(1)},\mathbf p_2^{(1)}}(\mathbf x_1),
 \overline\phi_{\mathbf p_1^{(0)},\mathbf p_2^{(0)}}(\mathbf x_0)
 \bigr).
\]
由坐标变换的构造和~\eqref{eq:a1-60-1}，立即得到
\begin{equation}
 \left\lvert
 \frac{\partial^{k+\ell}}{\partial T^k\partial\mathbf x^\ell}
 \left(
 \operatorname{ev}_{W_i}(w_{T;\mathbf x_1,\mathbf x_0}^{(1)\prime})
 -
 \operatorname{ev}_{W_i}(w_{T;\mathbf x'_1,\mathbf x'_0}^{(2)\prime})
 \right)
 \right\rvert
 <C_{k,\ell}\exp(-c_{k,\ell}T).
\tag{A1.60.2}\label{eq:a1-60-2}
\end{equation}
式~\eqref{eq:a1-60-1}、\eqref{eq:a1-60-2} 与
引理~\ref{lem:a1-58} 立即推出命题~\ref{prop:a1-56}。
\end{proof}

逐项检查第 7.1.3 小节的证明细节，就能直接证明引理~\ref{lem:a1-58}。
由于 $T$-导数的估计可能不那么标准，下面为完整起见详细证明该引理中的
指数估计。

证明之前，先回顾第 7.1.3 小节的证明方法仿照了 \textbf{[Fuk96II]}。
本书命题 7.1.27 与 \textbf{[Fuk96II]} 引理 8.5 中线性化算子满射性的证明，
都遵循唐纳森在 \textbf{[Don86I]} 中的“交替法”。第 7.1.3 小节的方法
略有不同，它使用隐函数定理。（不过，在该证明最重要的环节
引理~\ref{lem:7-1-32} 中，我们用的仍是交替法的同一思想。）
\textbf{[Fuk96II]} 用交替法研究方程的线性部分，而唐纳森在
\textbf{[Don86I]} 的原始论证中直接用交替法求解非线性偏微分方程；在他
所研究的情形中，这个方程就是反自对偶方程。

为透明地看出构造对 $T$ 的依赖，直接使用交替法，似乎比把证明的一部分
装进泛函分析的黑箱（即隐函数定理）更简单。自这种方法二十多年前首次
出现以来，证明所涉及的技巧如今都已成为标准。例如，唐纳森发表于 1986
年的论文 \textbf{[Don86I]}，已经成为严肃研究几何偏微分方程者的经典文献。
尽管如此，我们仍针对当前情形描述交替法的主要方案，再检查该方法所用
构造的细节，并推出引理~\ref{lem:a1-58}。特别地，我们将说明如何用这一
方法估计 $T$-导数。采用圆柱坐标时，$T$-导数的估计似乎比采用定义域节点
附近的标准坐标更容易。

\begin{proof}[引理~\ref{lem:a1-58} 的证明]
首先，可以假设 $(2n-2)$ 维子流形 $Y_0,Y_1$ 关于所取黎曼度量全测地。
在以下论证中，考虑满足约束
\[
 w(z_i^{\mathrm{int}})\in Y_i,
 \qquad i=0,1,
\]
的映射空间。因此，$w$ 处的切向量是
$(w^*TM,w|_{\partial\Sigma}^*TL)$ 的一个截面 $V$，满足
\[
 V(z_i^{\mathrm{int}})
 \in T_{w(z_i^{\mathrm{int}})}Y_i.
\]
这里 $\Sigma$ 可以是 $\Sigma_i$（$i=0,1$），也可以是 $\Sigma_T$。

从命题~\ref{prop:a1-56} 的证明中构造的近似解
$w_{T;\mathbf x_1,\mathbf x_0}^{(j)\prime}$ 开始。证明引理时可分别处理
每个 $j$，故从现在起省略 $j$，简写为
$w'_{T;\mathbf x_1,\mathbf x_0}$。

\medskip\noindent\textbf{（循环 0）}
置
\[
 w'_{T;\mathbf x_1,\mathbf x_0}
 =w'_{T;\mathbf x_1,\mathbf x_0;0}.
\]
循环 0 没有步骤 0-1 和步骤 0-2，因为我们已经从近似解开始；可与下述
较后循环中的相应步骤比较。

\smallskip\noindent\textbf{（步骤 0-3：对近似解施加
$\bar\partial$ 算子。）}
考虑 $\bar\partial$ 导数
\[
 \bar\partial(w'_{T;\mathbf x_1,\mathbf x_0;0})
 \in
 \Gamma\!\left(
 \Sigma_T;
 \Lambda^{0,1}\otimes
 (w'_{T;\mathbf x_1,\mathbf x_0;0})^*TM
 \right).
\]
按引理~\ref{lem:7-1-29} 的意义，它是指数小的。把它记为
\[
 \operatorname{Err}_0
 =\bar\partial(w'_{T;\mathbf x_1,\mathbf x_0;0}).
\]

\smallskip\noindent\textbf{（步骤 0-4：把误差项分成两片。）}
使用式~\eqref{eq:7-1-21} 下方定义的算子
$J^S_{*,i}$（$i=0,1$），得到
\[
 \operatorname{Err}_{0,i}
 =J^0_{*,i}(\operatorname{Err}_0)
 \in
 \Gamma\!\left(
 \Sigma_i;\Lambda^{0,1}\otimes
 (w'_{T;\mathbf x_1,\mathbf x_0;0;i})^*TM
 \right),
 \qquad i=0,1.
\]
这里在 $J^S_{*,i}$ 中取 $S=0$。下面解释
$w'_{T;\mathbf x_1,\mathbf x_0;0;i}$ 的记号。算子
$J^0_{*,i}$（$i=0,1$）用一组沿 $[-1,1]\times[0,1]$ 截断的单位分解，
把截面分成两片。

截断以后，截面 $J^0_{*,1}(\operatorname{Err}_0)$ 的支集位于
\[
 [5T+1,\infty)\times[0,1]\subset\Sigma_1
\]
之外。（注意，第 7.1.3 小节把 $\Sigma_1$ 记作
$X(t_1,\ell_1)$，并且它包含 $[0,\infty)\times[0,1]$。）定义映射
\[
 w'_{T;\mathbf x_1,\mathbf x_0;0;1}\colon\Sigma_1\longrightarrow M,
\]
使它在 $[10T-1,\infty)\times[0,1]$ 之外同已定义的
$w'_{T;\mathbf x_1,\mathbf x_0;0}$ 相同，并利用对应节点处的指数映射和
适当的单位分解，把它延拓到
$[10T-1,\infty)\times[0,1]$ 上。类似地定义
$w'_{T;\mathbf x_1,\mathbf x_0;0;0}\colon\Sigma_0\to M$。注意，
$[10T-1,\infty)\times[0,1]$ 这一部分在步骤 $n$-2 和步骤 $n$-4 中
不会使用；只要所取延拓的 $T$-导数等为指数小，怎样延拓都无关紧要。
这一点可对循环次数 $n$ 同时作归纳验证。

由于稍后在式~\eqref{eq:a1-62} 后说明的理由，我们在此结束第零个循环，
进入下一循环。

\medskip\noindent\textbf{（循环 1）}

\smallskip\noindent\textbf{（步骤 1-1：在两片上求解方程。）}
这里应用引理~\ref{lem:7-1-20} 的一个版本，即利用下列算子的满射性：
\begin{equation}
 \begin{aligned}
 &\pi\circ
 D_{w'_{T;\mathbf x_1,\mathbf x_0;0;1},
   w'_{T;\mathbf x_1,\mathbf x_0;0;0}}
 \bar\partial\colon\mathcal X\\
 &\hspace{4em}\longrightarrow
 \mathcal E^{0,p}_{\Sigma_1}/E_1
 \oplus
 \mathcal E^{0,p}_{\Sigma_0}/E_0.
 \end{aligned}
\tag{A1.61}\label{eq:a1-61}
\end{equation}
\translatornote{源书此处写“引理 29.20”，但紧接着连续把
式 (A1.61) 与引理 7.1.20 的算子比较；译文恢复为引理 7.1.20。
见疑误 A-0708。}
式~\eqref{eq:a1-61} 同引理~\ref{lem:7-1-20} 中出现的算子略有不同：
这里用 $w'_{T;\mathbf x_1,\mathbf x_0;0;i}$ 代替该引理中的 $w_i$。
其原因是，我们现在进行的是唐纳森交替论证的非线性版本，而引理
\ref{lem:7-1-20} 用于线性版本。不过，两者之差指数小，故
式~\eqref{eq:a1-61} 仍然满射。

把式~\eqref{eq:a1-61} 的一个右逆作用于
$(\operatorname{Err}_{0,1},\operatorname{Err}_{0,0})$，得到
\[
 V_{1,i}\in
 \Gamma\!\left(
 (\Sigma_i,z_i^{\mathrm{int}});
 ((w'_{T;\mathbf x_1,\mathbf x_0;0;i})^*TM,TY_i)
 \right),
 \qquad i=0,1.
\]

\smallskip\noindent\textbf{（步骤 1-2：胶合这些解。）}
像式~\eqref{eq:7-1-21} 下方构造 $w'_T$ 那样使用
$(V_{1,1},V_{1,0})$。更准确地说，按式~\eqref{eq:7-1-23} 胶合它们，
并在该式中取 $S=2T$，由此得到
\[
 \Gamma\!\left(
 (\Sigma_T,z_1^{\mathrm{int}},z_0^{\mathrm{int}});
 ((w'_{T;\mathbf x_1,\mathbf x_0;0})^*TM,TY_1,TY_0)
 \right)
\]
的一个元素。对它施加指数映射
$\exp_{w'_{T;\mathbf x_1,\mathbf x_0;0}}$，得到新映射
\[
 w'_{T;\mathbf x_1,\mathbf x_0;1}\colon\Sigma_T\longrightarrow M,
\]
并且
$w'_{T;\mathbf x_1,\mathbf x_0;1}(z_i^{\mathrm{int}})\in Y_i$
（$i=0,1$）。这里所用单位分解沿
\[
 [2T-1,2T+1]\times[0,1]\subset\Sigma_T
\]
截断；该区域对应于
$[7T-1,7T+1]\times[0,1]\subset\Sigma_0$。与步骤 $n$-4 中单位分解
所沿的截断区域相比，它朝端部方向远 $2T$。这 $2T$ 的差正是使误差项
相对于 $T$ 按指数阶下降的关键。唐纳森在 \textbf{[Don86I]} 中的论证，
以及引理~\ref{lem:7-1-32} 的证明，都以类似方式使用这一点。

\smallskip\noindent\textbf{（步骤 1-3：对近似解施加
$\bar\partial$ 算子。）}
与步骤 0-3 一样，考虑
\[
 \operatorname{Err}_1
 =\bar\partial(w'_{T;\mathbf x_1,\mathbf x_0;1})
 \in
 \Gamma\!\left(
 \Sigma_T;\Lambda^{0,1}\otimes
 (w'_{T;\mathbf x_1,\mathbf x_0;1})^*TM
 \right).
\]

\smallskip\noindent\textbf{（步骤 1-4：把误差项分成两片。）}
同步骤 0-4 一样使用 $J^S_{*,i}$（$i=0,1$），得到
\[
 \operatorname{Err}_{1,i}
 =J^0_{*,i}(\operatorname{Err}_1)
 \in
 \Gamma\!\left(
 \Sigma_i;\Lambda^{0,1}\otimes
 (w'_{T;\mathbf x_1,\mathbf x_0;1;i})^*TM
 \right),
 \qquad i=0,1.
\]
映射 $w'_{T;\mathbf x_1,\mathbf x_0;1;i}$ 的定义与步骤 0-4 相同。

现在又回到步骤 1-1 的同一情形，因而可以开始下一循环：
\[
 \text{（循环 2）}\qquad
 \text{（步骤 2-1）}\Longrightarrow
 \text{（步骤 2-2）}\Longrightarrow
 \text{（步骤 2-3）}\Longrightarrow
 \text{（步骤 2-4）},
\]
并以相同方式继续到循环 $n$。

与 \textbf{[Don86I]} 相同（再作 \textbf{[Fuk96]} 中适用于博特--莫尔斯
情形的小修改），可以证明极限
\[
 \lim_{n\to\infty}w'_{T;\mathbf x_1,\mathbf x_0;n}
\]
是伪全纯映射，且就是所需的
$w_{T;\mathbf x_1,\mathbf x_0}$。

注意，在每个循环 $n$ 开始时，我们有
\begin{equation}
 \Gamma\!\left(
 \Sigma_i;\Lambda^{0,1}\otimes
 (w'_{T;\mathbf x_1,\mathbf x_0;n-1;i})^*TM
 \right),
 \qquad i=0,1,
\tag{A1.62}\label{eq:a1-62}
\end{equation}
中的元素。定义域 $\Sigma_i$ 与 $T$ 无关；丛
$(w'_{T;\mathbf x_1,\mathbf x_0;n-1;i})^*TM$ 可能依赖于 $T$，但这种差别
相对于 $T$ 是指数小的，这也可以同时归纳证明。因此，可在靶空间 $M$
的局部取切丛平凡化，从而把这些丛彼此识别。也就是说，每个循环都从
式~\eqref{eq:a1-62} 中的元素开始，并在循环末尾得到同一函数空间中的
元素，故在每个循环中取 $T$-导数是有意义的。这正是我们按上述方式安排
每个循环起点的理由。

现在可以估计各步骤中出现的截面的 $T$-导数。不含 $T$-导数的各项，
由上面的描述很容易估计。

先考虑步骤 $n$-2、步骤 $n$-3 与步骤 $n$-4。这些步骤都是局部过程：
用单位分解胶合截面、施加指数映射，以及施加 $\bar\partial$ 算子。这些
过程本身都不含参数 $T$，但过程中出现的一些函数等含有 $T$，因而成为
非零 $T$-导数的来源。检查构造可知，这类非零导数有如下两个来源：

\smallskip\noindent\textup{(A)}
我们使用依赖于 $T$ 的单位分解。对含单位分解的项取 $T$-导数，会产生
形如 $d^k\chi_T/dT^k$ 的项。这些项的 $C^\infty$ 范数有界，故乘以它们
是有界算子。另一方面，它们所乘的截面，其指数衰减已经在归纳的较早
阶段建立。因此，这些项可以按指数阶估计。

\smallskip\noindent\textup{(B)}
构造过程中会出现映射
$w'_{T;\mathbf x_1,\mathbf x_0;n}$ 或
$w'_{T;\mathbf x_1,\mathbf x_0;n;i}$（$i=0,1$）。它们依赖于 $T$，因为
$T$ 已在归纳的较早阶段出现。不过，可以归纳证明：由这种 $T$ 依赖产生的
变分，在这里所用的局部加权 $W^{1,p}$ 范数下是指数小的（该范数的定义
见第 7.1.3 小节）。因此，取 $T$-导数会得到“乘以某个截面”的算子，
而该截面的局部 $W^{1,p}$ 范数为指数小量级。再利用乘法映射
\[
 W^{1,p}\times W^{1,p}\longrightarrow L^p
\]
的有界性，便知由此产生的项也指数小。
\translatornote{源书字面上说上述映射本身在局部加权 $W^{1,p}$ 范数下
指数小；但它们取值于 $M$，本段实际估计的是其 $T$ 依赖所产生的变分。
译文按紧前“非零 $T$-导数的来源”及紧后乘法截面的型别补明。见疑误
A-0712。}

再考虑步骤 $n$-1。这个步骤稍微更不平凡，因为它是非局部的，并且使用
式~\eqref{eq:a1-61} 的右逆。为研究此步骤，考察该式对 $T$ 的导数。
式~\eqref{eq:a1-61} 的首项是柯西--黎曼算子，不含 $T$；它的零阶部分则
含 $T$，因为映射 $w'_{T;\mathbf x_1,\mathbf x_0;n;i}$ 含 $T$。与上述
(B) 相同，式~\eqref{eq:a1-61} 的 $T$-导数是“乘以某个截面”的算子，
而该截面的加权 $W^{1,p}$ 范数指数小。还要注意，在步骤 $n$-1 中，
式~\eqref{eq:a1-61} 的函数空间完全不依赖于 $T$。因此，利用标准论证
（例如诺伊曼级数），可以找到它的右逆，使该右逆的 $T$-导数，相对于
$L^p\to W^{1,p}$ 算子的算子范数指数小。

这就完成了引理~\ref{lem:a1-58} 的证明。
\end{proof}

\begin{remark}{A1.63}\label{rem:a1-63}
若把胶合圆盘的双倍覆盖识别为一个闭黎曼曲面，则在奇点附近，胶合曲面可
局部看成方程
\[
 z_1z_2=r,
 \qquad r=\exp(-cT)
\]
的解。因此，(A1.57.3) 可以写成
\[
 \left\lvert
 \frac{\partial^k
  (\phi'_{\mathbf p_1,\mathbf p_2}
   -\phi_{\mathbf p_1,\mathbf p_2})}
 {\partial r^k}
 \right\rvert
 <r^{c_k-k}.
\]
\translatornote{源书左端只写
$|\partial^k\phi_{\mathbf p_1,\mathbf p_2}/\partial r^k|$；但它声称这是
(A1.57.3) 的改写，而后者估计的是 $\phi'-\phi$。特别在 $k=0$ 时，
源式会与 (A1.57.2) 的非零极限冲突，故译文补回差值。见疑误 A-0709。}
这与 $f(r)=r^\lambda$ 一类函数所满足的估计同型；在这个意义下，我们的
估计很自然。另一方面，它还不足以推出
$\phi_{\mathbf p_1,\mathbf p_2}$ 关于 $r$（或 $r$ 的某个幂）的光滑性。
这里选取的坐标是
\[
 r=e^{-1/s}.
\]
在这个坐标中，$f(s)=e^{-\lambda/s}$ 当然光滑。

或许可以期待 $\phi_{\mathbf p_1,\mathbf p_2}$ 关于 $r$ 的某个幂光滑。
例如，当可以在实代数几何范畴中工作时，似乎就是如此。这里不尝试改进
估计，因为我们的目的不需要更强的估计，尽管这种估计本身看起来很有
意思。作者最近（在 2007 年）得知，有几位研究者正在尝试证明这一更强
估计。
\end{remark}

注~\ref{rem:a1-63} 的一个推论是：若在
$\mathcal M_{k+1}^{\mathrm{main}}$ 上赋予标准 $C^\infty$ 结构，则遗忘映射\liftIndex{I389}
\[
 \mathcal M_{k+1}^{\mathrm{main}}(L;\beta)
 \longrightarrow
 \mathcal M_{k+1}^{\mathrm{main}}
\]
在源的仓西邻域上，关于我们的坐标并不光滑。我们从未使用这个映射的
光滑性。不过，这同映射~\eqref{eq:7-3-7-2} 的光滑性有关，即
\[
 \operatorname{forget}_{\vec a}\colon
 \mathcal M_{k+1}^{\mathrm{main}}(\beta;\vec P^+)
 \longrightarrow
 \mathcal M_{k+1-|\vec a|}^{\mathrm{main}}(\beta;\vec P).
\]
（这个问题在本书中只出现于同伦幺元的构造。）

下面说明这一问题及其解决办法。为简化记号，考虑如下情形。令
\[
 \Sigma=D_1^2\cup D_2^2\cup D_3^2,
 \qquad
 \{p_i\}=D_i^2\cap D_{i+1}^2\ (i=1,2),
 \qquad D_1^2\cap D_3^2=\varnothing.
\]
设
\[
 w\colon(\Sigma,\partial\Sigma)\longrightarrow(X,L)
\]
是伪全纯映射，置 $w_i=w|_{D_i^2}$、$\beta_i=[w_i]$。假设 $w_2$ 是
常值映射，因而 $\beta_2=0$。取
$z_0\in D_1^2,z_1\in D_2^2$，并置
$\beta=\beta_1+\beta_3$。于是
\[
 [(\Sigma,w,(z_0,z_1))]\in\mathcal M_2^{\mathrm{main}}(\beta).
\]
遗忘 $z_1$ 后，有
\[
 \operatorname{forget}([(\Sigma,w,(z_0,z_1))])
 =[(\Sigma',w',(z_0))]
 \in\mathcal M_1^{\mathrm{main}}(\beta),
\]
其中 $\Sigma'$ 由 $D_1^2,D_3^2$ 在一点胶合而成，$w'$ 则以显然方式由
$w$ 得到。

下面描述遗忘映射在仓西邻域中的形状。考虑
\[
 \begin{aligned}
 &[(D_1^2,w_1,(z_0,p_1))]
   \in\mathcal M_2^{\mathrm{main}}(\beta_1),\\
 &[(D_2^2,w_2,(p_1,z_1,p_2))]
   \in\mathcal M_3^{\mathrm{main}}(0),\\
 &[(D_3^2,w_3,(p_2))]
   \in\mathcal M_1^{\mathrm{main}}(\beta_3).
 \end{aligned}
\]
令 $U_i$（$i=1,2,3$）为它们的仓西邻域，$E_i$ 为相应障碍丛。
$[(\Sigma,w,(z_0,z_1))]$ 的仓西邻域是
\[
 (U_1\mathbin{\times_L}U_2\mathbin{\times_L}U_3)
 \times[0,\epsilon)^2,
\]
而 $[(\Sigma',w',(z_0))]$ 的仓西邻域是
\[
 (U_1\mathbin{\times_L}U_3)\times[0,\epsilon).
\]
\translatornote{源书把第三分量印成
$[(D_1^3,w_3,(z_2)]\in\mathcal M_1^{\mathrm{main}}(\beta_2)$，又把目标邻域
印成 $(U_1\times_L\times U_3)\times[0,\epsilon)$。依本段已经定义的
三个圆盘、两个节点、$\beta_2=0$ 与 $\beta=\beta_1+\beta_3$，译文分别
恢复为 $(D_3^2,p_2,\beta_3)$ 和 $U_1\times_LU_3$。见疑误
A-0710、A-0711。}
这里所关心的是分层法向方向上的坐标，即
$[0,\epsilon)^2$ 与 $[0,\epsilon)$ 的坐标。下面依本小节的讨论回顾这些
坐标。

对 $[(\Sigma,w,(z_0,z_1))]$，有两个奇点，故有两个参数 $T_1,T_2$；
它们分别对应于圆柱坐标中两个颈的长度。我们为 $[0,\epsilon)^2$ 选取的
坐标是
\[
 s_1=1/T_1,
 \qquad s_2=1/T_2.
\]
对 $[(\Sigma',w',(z_0))]$，只需解消一个奇点。令 $T$ 为相应参数，即该颈
的长度，而 $[0,\epsilon)$ 的坐标是 $s=1/T$。遗忘 $z_1$ 后，有
\[
 T=T_1+T_2+c,
\]
其中 $c$ 是
$(U_1\times_LU_3)\times[0,\epsilon)$ 上的某个函数；当
$T_1,T_2\to\infty$ 时，它连同全部导数都按指数阶趋于一个常数。于是
\[
 s=\frac1{T_1+T_2+c}
 =\frac{s_1s_2}{cs_1s_2+s_1+s_2}.
\]
\translatornote{源书第一项的分母印作
$T_1^{-1}+T_2^{-1}+c$；但紧前已经给出 $T=T_1+T_2+c$ 与 $s=1/T$，
而同一行右端也正是 $1/(T_1+T_2+c)$ 用 $s_i=1/T_i$ 改写后的结果。
故译文恢复为 $T_1+T_2+c$。见疑误 A-0713。}
它在 $s_1=s_2=0$ 处不光滑，因而通过
$\operatorname{forget}$ 拉回多值截面时会遇到困难。

解决办法如下。注意，我们对扰动（多值截面）的构造是逐层构造：从较低
维层开始，把它延拓到该层的邻域，再利用横截性定理继续延拓。关键在于，
若多值截面在某一层上横截，那么向其邻域的任意延拓仍然横截。因此，对
这种延拓有很大的选择自由。

利用这一自由，把截面延拓到该层的邻域，使其沿 $T$ 坐标按指数阶衰减。
在当前例子中，要求对 $k\geq1$，多值截面 $\mathfrak s$ 的每个分支都满足
\[
 \left\lvert
 \frac{\partial^{k+\ell}\mathfrak s}
      {\partial T^k\partial\mathbf x^\ell}
 \right\rvert
 <C_{k,\ell}\exp(-c_{k,\ell}T).
\]
引理~\ref{lem:a1-58} 蕴含：这一条件在所得仓西结构的坐标变换下不变。
显然，满足该假设的多值截面经上述拉回后仍满足同一条件；特别地，这个
拉回是光滑的。因此，整个构造都可以在满足上述条件的多值截面中进行。
